\documentclass[10pt]{article}

\usepackage[margin=0.85in]{geometry}
\usepackage[T1]{fontenc}
\usepackage[utf8]{inputenc}
\usepackage{lmodern}
\setlength{\parindent}{0pt}
\setlength{\parskip}{4pt}
\pagestyle{empty}
\newcommand{\proposalsection}[1]{\vspace{2pt}\textbf{\normalsize #1}\par}

\begin{document}

\begin{center}
{\large \textbf{LLM Hallucination Detection via Uncertainty}}\\
Mustafacan Koc\\
SEDS 537 -- Machine Learning\\
Spring 2026
\end{center}

\proposalsection{Project Title}
LLM Hallucination Detection via Uncertainty-Aware Multi-Signal Fusion.

\proposalsection{Problem Description}
Large language models (LLMs) can generate fluent but factually incorrect or unsupported responses, known as hallucinations. These errors reduce the reliability of LLM-based systems in tasks where factual accuracy is important. This project aims to detect hallucinated outputs using uncertainty signals derived from the generation process, such as token probabilities and entropy. The task will be formulated as a binary classification problem in which a generated response is classified as supported or hallucinated.

\proposalsection{Proposed Approach}
The proposed method is a multi-signal hallucination detector that combines token-level uncertainty features, response consistency, and retrieval-based evidence agreement. Uncertainty features such as mean entropy, maximum entropy, average token confidence, and low-confidence token ratio will be extracted from generated answers. In addition, multiple responses to the same prompt can be compared to measure self-consistency, while a lightweight RAG component will retrieve relevant passages and estimate whether the answer is supported by external evidence. These signals will be combined in a classifier to predict whether a generated response is hallucinated.

\proposalsection{Baseline Methods}
The baseline methods will include entropy thresholding, token confidence thresholding, Logistic Regression, and Support Vector Machine (SVM). The threshold-based methods will classify responses using simple cutoffs on uncertainty signals such as token entropy and token probabilities. Logistic Regression and SVM will use handcrafted uncertainty features, including mean entropy, maximum entropy, average token confidence, and low-confidence token ratio, to distinguish hallucinated and supported responses.

\proposalsection{Dataset}
The main datasets will be TruthfulQA and HaluEval. TruthfulQA provides question-answering prompts designed to evaluate factual truthfulness, while HaluEval contains examples labeled for hallucination behavior. One dataset can be used for development and threshold tuning, and the other can be used for external validation to test generalization across benchmarks.

\proposalsection{Evaluation Strategy}
The methods will be evaluated as binary hallucination detectors using accuracy, precision, recall, F1-score, AUROC, and PR-AUC. The proposed method will be compared against all baselines on both datasets. An ablation study will measure the impact of uncertainty features, self-consistency, and retrieval-based evidence signals, and an error analysis will examine overconfident failures and unsupported but fluent responses.

\proposalsection{Expected Challenges}
Expected challenges include poorly calibrated token probabilities, hallucinated answers that appear fluent and confident, retrieval failures caused by weak evidence, and differences in label definitions across datasets. In addition, collecting multiple generations for consistency analysis may increase computational cost.

\end{document}
